\scanpage{16}

\deleted{ouverte, et $\bar f$ propre. Ex] morphisme
quasi-projectif sur $Y$ loc. noeth.}

13) \uline{Faisceaux} \uline{constructibles}

Faisceaux loc. \deleted{isom.} constants

Faisceaux constructibles : en général, il est
loc. constant : fibres \deleted{finies} de t.f.
sur des ens loc fermés cons. qui partitionnent
$\unsure{U|X}$. Dans le cas noethérien, critère
par adhérences \unsure{d'un} \unsure{pt} \deleted{\unsure{générique}}.

\uline{Th.}\quad Si $f:X\longrightarrow Y$ \uline{propre}, $Y$ \uline{loc noeth},
alors pour $\unsure{\forall|\text{tt}}$ faisceau \uline{de torsion}
\uline{constructible} $F$ sur $X$, les $R^i f_*(F)$
sur $Y$ sont constr.

\uline{Cor.}\quad $X$ propre sur un corps sép. clos
$k$. Alors pour $\unsure{\forall|\text{tt}}$ $F$ de torsion
constructible sur $X$, les $H^i(X,F)$ sont
des groupes finis.

\uline{Remarque}\quad \unsure{Question} de la généralisation

\mspage{15}
